% ============================================================================
%  A — file phần: Nền tảng nghĩa
%  Ngày tạo: 2026-09-06 | Sửa gần nhất: 2026-09-06 | Ôn gần nhất: chưa
% ============================================================================
\documentclass[../english.tex]{subfiles}
\begin{document}
\part{NỀN TẢNG NGHĨA}
\begin{tomtat}
Phần A dựng bốn thứ mà mọi chương sau đều dùng: một cách hiểu đúng về việc nghĩa nằm ở đâu (chương 1), một danh sách rõ ràng về việc ``biết một từ'' gồm những gì và cần bao nhiêu từ để đọc được tài liệu chuyên ngành (chương 2), lượng ngữ pháp tối thiểu để tách một câu dài thành các mệnh đề (chương 3), và cách đọc cấu tạo từ để đoán nghĩa từ mới mà không cần tra (chương 4). Đọc xong phần này, bạn thay được thói quen ``tra từ, ghi một nghĩa tiếng Việt'' bằng một cách ghi nhớ giữ được đủ thông tin để \emph{dùng lại} từ đó. Nếu chỉ nhớ một điều: nghĩa của một từ được quyết định bởi những từ thường đứng cạnh nó, nên học từ mà không học ngữ cảnh là học một nửa.
\end{tomtat}

Có một lý do rất thực dụng để phần này đứng trước. Người học tiếng Anh lâu năm thường có vốn từ lớn hơn họ tưởng, nhưng vốn đó ở dạng \emph{thụ động và nghèo thông tin}: mỗi từ gắn với đúng một nghĩa tiếng Việt, không kèm ngữ cảnh, không kèm sắc thái, không biết nó đi với giới từ nào. Vốn từ như vậy đủ để đoán ý chung của một đoạn văn, nhưng không đủ để phân biệt hai câu chỉ khác nhau một từ --- mà trong tài liệu khoa học, khác biệt đó thường là toàn bộ nội dung.

Chương 1 và 2 xử lý đúng chỗ đó. Chương 1 giải thích vì sao từ điển không bao giờ đủ và nghĩa thực sự được quyết định thế nào; chương 2 chuyển nhận xét ấy thành một danh sách kiểm cụ thể cho việc học từ, kèm các con số về ngưỡng từ vựng cần cho việc đọc.

Chương 3 và 4 là hai công cụ. Ngữ pháp ở đây được trình bày ở mức tối thiểu \emph{có chủ ý}: chỉ đủ để bạn nhìn một câu bốn dòng trong một bài báo và tách được đâu là mệnh đề chính, đâu là bổ nghĩa, cái gì bổ nghĩa cho cái gì. Đó là mức ngữ pháp thật sự cần cho việc đọc, và nó nhỏ hơn nhiều so với chương trình ngữ pháp thông thường. Chương 4 cho công cụ thứ hai: khi gặp từ lạ trong tài liệu chuyên ngành, phần lớn thời gian bạn đoán được nghĩa từ gốc và phụ tố, và việc đoán được ấy quan trọng hơn việc tra, vì nó giữ cho mạch đọc không đứt.
\subfile{00-map}
\subfile{01-nghia-nam-o-dau/nghia-nam-o-dau}
\subfile{02-biet-mot-tu/biet-mot-tu}
\subfile{03-ngu-phap-toi-thieu/ngu-phap-toi-thieu}
\subfile{04-tu-nguyen-va-cau-tao/tu-nguyen-va-cau-tao}
\ifSubfilesClassLoaded{\printbibliography[title={Nguồn}]}{}
\end{document}
